%=====================================================================
\needspace{16\baselineskip}\hdr{D}{附录 D：报告速查}

\subhdr{D.1 报告要交什么}

\begin{center}
\begin{minipage}{\linewidth}
\begin{tabular}{@{}p{3.6cm}p{3.2cm}p{7.0cm}@{}}
\toprule
\thead{项} & \thead{什么时候} & \thead{内容} \\
\midrule
预习 8 题 & 进实验室前 & 英文答案写进报告（下面 D.2 可直接抄） \\
数据表 & 课上 & 打印 Preview p3 或本书 J1/J2/J3 的表 \\
数据处理 & 课后 & 六项：$\alpha$、$\delta_{\min}$、$n$、$\varphi_k$、$d$、$D$（J1--J4 的 \S ⑤） \\
口试 & 课上可能被问 & 12 问要点（D.3），不必写进报告 \\
选做 & 自愿 & 3 道扩展题（D.4） \\
\bottomrule
\end{tabular}
\end{minipage}
\end{center}

\subhdr{D.2 预习 8 题：中英答案}

每题先给\textbf{可直接抄进报告的英文}（压到 50--60 词），再给\textbf{中文解释}（理解用，不必抄）。

\subhdr{Q1\quad What are the requirements for fine adjustment of the spectrometer?}

\textbf{EN.} Three conditions: (1) the collimator emits a parallel beam, i.e. its slit lies in the focal plane of the collimator lens; (2) the telescope receives parallel light, i.e. it is focused at infinity with the reticle in the focal plane of the objective; (3) both optical axes are perpendicular to the rotation axis, and the optical surface of the element under test is parallel to it.

\textbf{中文.} 三条：准直管出平行光、望远镜收平行光、两轴都垂直主轴且待测元件的光学面平行主轴。第三条最容易只做一半——前两条只保证「有平行光」，第三条才保证「角度读数有意义」。

\subhdr{Q2\quad What is the purpose of visual coarse adjustment? How to carry it out?}

\textbf{EN.} Purpose: to bring both axes roughly perpendicular to the rotation axis and the stage roughly horizontal, so that the autocollimated image can enter the small field of view; the fine screws move only a few degrees. Method: by eye set the telescope (No.\,12) and the collimator (No.\,27) at the same height facing each other, then level the stage with its three screws.

\textbf{中文.} 目的是把三条轴摆进「细调能起作用」的范围。讲义原文说粗调「是细调成功的关键」，不是客套。

\subhdr{Q3\quad Explain the principle of focusing on infinity by the self-aligning method.}

\textbf{EN.} The illuminated transparent cross acts as the object. If the reticle lies in the focal plane, its light leaves the objective parallel and, after reflection by a plane mirror, returns parallel and is focused back onto the reticle plane; image and crosshair then show no parallax. Eliminating parallax therefore focuses the telescope at infinity.

\textbf{中文.} 关键是一条链：物在焦面 $\Leftrightarrow$ 出射平行 $\Leftrightarrow$ 返回后仍聚焦在焦面 $\Leftrightarrow$ 无视差。把这条链说完整，比只背「调到这里没有视差」得分高。

\subhdr{Q4\quad Draw a sketch of the plane mirror placement; how to judge the axis is perpendicular?}

\textbf{EN.} (Sketch: the stage as a circle with three screws; the mirror standing vertically, one face toward the telescope.) Judgment: the green cross from the front face coincides with the upper crosshair; after turning the stage by $180^{\circ}$ the image from the rear face must also coincide. Both coinciding means the telescope axis is perpendicular to the rotation axis.

\textbf{中文.} 画图要素：三个螺钉、镜子竖直、一个反射面朝望远镜。判据是\textbf{两面都重合}，不是一面。

\subhdr{Q5\quad How to judge that the parallel light from the collimator is perpendicular to the rotation axis?}

\textbf{EN.} Align the telescope with the collimator and watch the slit image. Turn the slit sleeve until the image is parallel to the horizontal crosshair, then adjust screw No.\,27 until the image is centred vertically; the collimator axis is then perpendicular to the rotation axis. Finally turn the slit by $90^{\circ}$ and lock it.

\textbf{中文.} 判据两层：像\textbf{居中}（轴垂直主轴）与像\textbf{竖直}（狭缝本身竖直，这是后面读衍射角的前提）。

\subhdr{Q6\quad Sketch the reflection method for the vertex angle, and give the formula.}

\textbf{EN.} (Sketch: telescope at I for the beam from face $AB$, at II for the beam from $AC$.) Rotating a mirror by $\gamma$ turns the reflected ray by $2\gamma$; the two faces differ by $\alpha$, so the two beams differ by $2\alpha$ and the telescope turns through $2\alpha$. Hence
\[
\alpha=\frac14\bigl(|\phi_2-\phi_1|+|\phi_2'-\phi_1'|\bigr),
\]
with $\phi_1,\phi_1'$ and $\phi_2,\phi_2'$ the vernier readings at I and II.

\textbf{中文.} 图要画出准直管、顶点、两个反射面、两束反射光、望远镜两个位置与四个角标。$\tfrac14$ 来自「双面镜的 $\tfrac12$」乘「双游标的 $\tfrac12$」。

\subhdr{Q7\quad Sketch the method of measuring the minimum deflection angle.}

\textbf{EN.} (Sketch: the prism in both orientations, the two emergent directions on either side of the fixed incident direction, $2\delta_{\min}$ apart.) Rotate the stage slowly: the slit image moves, stops, then reverses; the stationary point is the minimum deviation. Read both verniers, turn the prism to the other face and repeat. The difference is $2\delta_{\min}$, so
\[
\delta_{\min}=\frac14\bigl(|\phi_2-\phi_1|+|\phi_2'-\phi_1'|\bigr).
\]

\textbf{中文.} 图要能看出「入射方向不变、出射方向在两侧各偏 $\delta_{\min}$」。文字部分必须写「像停住后折返」这个判据——那是操作的核心。

\subhdr{Q8\quad Explain the principle of measuring grating constants; give the grating equation.}

\textbf{EN.} For normal incidence, light from two adjacent slits has a path difference $d\sin\varphi$ in the direction $\varphi$; a bright line appears when this equals an integer number of wavelengths:
\[
d\sin\varphi=k\lambda\qquad(k=0,\pm1,\pm2,\dots).
\]
Measuring $\varphi_1$ for a known $\lambda$ gives $d=\lambda/\sin\varphi_1$, and the line density is $N=1/d$ ($d$ in mm).

\textbf{中文.} 答题必须把「光程差 $d\sin\varphi$」这一步写出来，只写公式不写来源通常扣分。

\needspace{20\baselineskip}\subhdr{D.3 口试 12 问：要点}

\begin{center}
\begin{minipage}{\linewidth}
\small
\begin{tabular}{@{}p{0.5cm}p{4.8cm}p{8.6cm}@{}}
\toprule
\thead{\#} & \thead{问题} & \thead{要点} \\
\midrule
1 & 分光计由几部分组成？各起什么作用？ & 五部分：望远镜（收平行光、自准直定位方向）、准直管（出平行光，把缝的位置换成方向）、载物台（放待测件、绕主轴转）、度盘与游标（把方位读成 $1'$）、底座（支撑与定位主轴） \\
2 & 主截面是什么？ & 垂直于棱镜两折射面交线的截面。只有主截面内的光满足 $\alpha=r+r'$ 与 $\delta=i+i'-\alpha$；主截面垂直主轴是全部角度公式的前提 \\
3 & 为什么要求望远镜光轴平行于主截面？怎么调？ & 因为公式在主截面内推。调法是让主截面垂直主轴：$AB$、$AC$ 两面自准直像都与上叉丝重合，只调 $a_1,a_3$ \\
4 & 自准直的步骤与「能接收平行光」的判据？ & 四步：看清叉丝 → 找绿十字 → 移目镜筒至清晰 → 消视差。判据「清晰且无视差」，因为无视差 $\Leftrightarrow$ 分划板在焦面 $\Leftrightarrow$ 出射与返回都平行 \\
5 & 为什么绿十字与\textbf{上叉丝}重合？ & 透光「十」在轴下 $a$、上叉丝在轴上 $a$，两者对称；出射平行光与轴成 $\theta=a/f$，反射只反转轴向分量，返回光向另一侧偏，落在 $+a$ \\
6 & 细调的要求有哪些？ & 同 Q1 三条：准直管出平行光、望远镜聚焦无穷远、两轴垂直主轴 \\
7 & 找不到平面镜的「$+$」像怎么办？ & 回目测粗调：查电源与绿 LED、镜子是否放平且反射面朝望远镜、12/13 号是否偏太多、9 号镜筒是否把像推出视场 \\
8 & 测顶角时找不到棱镜反射像怎么办？ & 先回粗调；再查棱镜方位（$AB\perp a_1a_2$ 等）、顶点是否朝准直管、台面高矮是否挡光 \\
9 & 怎样判断最小偏向角？ & 慢转台并跟踪像，像\textbf{停住后折返}处即 $\delta_{\min}$；依据是 $\delta(i)$ 在 $i=i'$ 处取极小 \\
10 & 入射光不垂直光栅面会怎样？ & 起作用的缝距变成 $d\cos\gamma$，测出的 $d$ 偏大；$\pm k$ 不再对称、谱线不等高。必须用 $b_1,b_2$ 调垂直 \\
11 & 光栅光谱与棱镜光谱有何不同？哪条线偏折最大？ & 光栅：多级、近似等间距、长波偏折大（红在外）；棱镜：单套、不等间距、短波偏折大（紫在外）。光栅里是 $579.0$ nm 黄线，棱镜里是 $404.7$ nm 紫线 \\
12 & 缝太宽或太窄会怎样？ & 太宽：像成宽带，对准中心带主观误差，两近线并成一条；太窄：太暗，衍射使边缘发虚。取 $0.5$ mm 折中 \\
\bottomrule
\end{tabular}
\end{minipage}
\end{center}

\textbf{两问补充}：\textbf{（讲义 Q7）两条谱线不等高}——竖丝对准斜谱线时，交点位置随高度变，$\theta_k$ 带误差，且 $\pm k$ 不对称，所以要 $b_3$ 调等高。\textbf{（预习 III.9）注意事项}——不碰光栅与棱镜的光学面（指纹擦不掉）、汞灯不频繁开关且不直视、不自己合闸、读数前确认 16 号没动、狭缝与光栅调好后不再碰。

\subhdr{D.4 扩展 3 题（选做）}

\textbf{1. 除反射法外怎么测顶角：自准直法。}用自准直分别让望远镜对准 $AB$ 面与 $AC$ 面的\textbf{法线}（绿十字与上叉丝重合），两次位置之差是\textbf{两法线夹角}；两外法线夹角为 $180^{\circ}-\alpha$，故 $\alpha=180^{\circ}-$ 转角。

\warnbox{\textbf{一个巧合会骗人}：等边棱镜（$\alpha=60^{\circ}$）下，反射法转 $2\alpha=120^{\circ}$，自准直法转 $180^{\circ}-\alpha=120^{\circ}$——\textbf{两种方法转的角一样}。所以只看「转了 $120^{\circ}$」分辨不出用了哪种方法；换 $\alpha=45^{\circ}$ 的棱镜，反射法转 $90^{\circ}$、自准直法转 $135^{\circ}$，差别才出来。作答要写公式，不要只写数字。}

\textbf{2. 除最小偏向角法外怎么测折射率：掠入射法（极限角法）。}让光接近 $90^{\circ}$ 掠射入 $AB$ 面，从 $AC$ 面射出的光有一个\textbf{极限角} $\delta$（再转就消失）。测出它：

\[
n=\sqrt{1+\left(\frac{\sin\delta+\cos\alpha}{\sin\alpha}\right)^{2}}.
\]

推导：掠入射时 $i\to90^{\circ}$，故 $\sin r=1/n$；由 $r+r'=\alpha$ 得 $r'=\alpha-r$；再展开 $n\sin r'=\sin\delta$：

\[
\sin\delta=n\sin(\alpha-r)=n(\sin\alpha\cos r-\cos\alpha\sin r)=\sin\alpha\sqrt{n^{2}-1}-\cos\alpha,
\]

移项平方即得上式。核对：$n=1.5$、$\alpha=60^{\circ}$ 时 $\sin\delta=0.4682$、$\delta=27.92^{\circ}$，代回得 $n=1.500$ $\checkmark$。这个方法不需要找折返点，只要找出射光消失的边界，也是阿贝折射仪的原理基础。

\textbf{3. 分光计还能测什么}：未知波长（已知 $d$ 测 $\varphi_k$，$\lambda=d\sin\varphi_k/k$）；光栅色分辨本领（钠双线能否分开，$R=\lambda/\Delta\lambda$）；棱镜色散曲线（多条谱线的 $\delta_{\min}$ 得 $n(\lambda)$）；棱镜角色散（$D=\Delta\delta/\Delta\lambda$）；多种介质的折射率（最小偏向角法、掠入射法、临界角法）；谱线精细结构（高线数光栅高级次）。

\subhdr{D.5 出门清单}

\begin{itemize}
\item \textbf{预习 3 页的打印件}（p1 题目、p2 思考题、\textbf{p3 数据表}——原文写明「Print it yourself and bring it to class」）。
\item 本书 J1/J2/J3 的表（若不用预习表，就印这一份，别带两张）。
\item 铅笔、橡皮、直尺（作图题要用）、计算器。
\item 校园卡/学生证（签到、借仪器可能要用）。
\end{itemize}

\vspace{6pt}
\begin{center}
{\color{petal}\small\ornstar}\hspace{6pt}{\small\color{inkgray}全书完：J0 调好 $\to$ J1 顶角 $\to$ J2 折射率；J3 衍射角 $\to$ J4 光栅常量与角色散}\hspace{6pt}{\color{petal}\small\ornstar}
\end{center}
